\subsection*{S1.21 Exact clinical and tumor annotations of existing variants}

\paragraph{Population and exact matching.} All 4,669 unique GRCh38 REF/ALT variants from the existing 7,262-event ASB population were fixed before current database matching. The full archive contains 5,409 training, 1,457 validation and 396 chromosome-8/15 events; all labels and predictions were previously examined. Annotation inclusion requires identical chromosome, one-based position, genomic-plus REF and ALT. Clinical and tumor sources annotate this population; they do not change its binding-direction eligibility or supply a new binding measurement.

The official dated ClinVar GRCh38 file \path{clinvar_20261011.vcf.gz}, posted on 12 October and retrieved on 14 October 2026, was frozen rather than using the mutable latest alias. Its compressed 200,437,060 bytes passed the official MD5 check; source SHA-256, header, retained original VCF rows and retrieval failures are archived. Scanning 4,594,764 VCF records gave 520 exact matches and 4,149 unmatched variants. The match groups comprise 453 benign/likely benign, three P/LP, four uncertain significance, seven conflicting classifications, one other classification and 52 without germline classification. Missing records and missing germline classifications remain distinct. The VCF covers precisely mapped simple alleles, rather than every ClinVar genotype or assertion. Germline significance, oncogenicity and somatic clinical impact are kept in their respective fields \citep{landrum2018clinvar}. Definitions and source: \url{https://www.ncbi.nlm.nih.gov/clinvar/docs/ftp_primer/} and \url{https://www.ncbi.nlm.nih.gov/clinvar/docs/clinsig/}.

\paragraph{Direct clinical records and condition boundaries.} Direct NCBI ESummary records matched canonical SPDI exactly for all three P/LP sites: ASL at \path{chr7:66081825} G$>$A, record \path{VCV000092360.115}; ELP1 at \path{chr9:108880051} C$>$T, record \path{VCV001066986.10}; and ASXL1 at \path{chr20:32434485} C$>$A, record \path{VCV002498883.28}. Each has two-star germline review, criteria provided by multiple submitters without conflicts. ELP1 is aggregate likely pathogenic; ASL and ASXL1 have aggregate pathogenic/likely pathogenic classifications. The ELP1 VCV reports familial dysautonomia, medulloblastoma and not-provided conditions, but the respective RCV condition-level classifications could not be independently retrieved; no per-condition pathogenicity or cancer-driver conclusion is made. ASXL1's reported condition is only not provided. All direct responses, accessions, review status, last-evaluated dates and failures are retained.

\paragraph{Open TCGA occurrence.} The public GDC API was queried for all 4,669 keys in 19 batches with zero unresolved queries. Source status reported API 8.5.0 and GDC data release 46.0 (10 August 2026). Only TCGA project occurrences were retained. Six exact GRCh38 matches have one distinct case occurrence each: EIF3F in LUSC, HSP90B1 in UCEC, two MAP2K3 alleles in BRCA/STAD, C3 in LGG and ASXL1 in BLCA. All 57 transcript annotations are retained, separating canonical and alternate-transcript consequences. A coding mutation or an alternate-transcript UTR annotation does not establish an RNA-binding disease mechanism. Successful source absence is not biological absence. The archived project denominators (504, 559, 1,098, 443, 516 and 412, respectively) count cases with SSM records, not cases callable at each locus; any retained ratio is explicitly an indexed-case fraction rather than mutation prevalence. No controlled-access patient data were requested \citep{grossman2016gdc}. Endpoint definitions: \url{https://docs.gdc.cancer.gov/API/Users_Guide/Appendix_A_Available_Fields/}.

\paragraph{Complete joined display and case selection.} Let $C$ denote all three clinical P/LP variants and $T$ all six exact TCGA matches. The displayed union is $U=C\cup T$, with $|U|=3+6-1=8$ because ASXL1 overlaps. Every corresponding ASB event is retained: six clinical and ten TCGA events, with the three ASXL1 events counted once in the 13-event union. This object has 11 training and two validation events and no chromosome-8/15 event. Events are variant--RBP--cell observations, not independent patients. Variant rows follow genomic order; no model response determines inclusion or display order. ELP1 is the sole splice-donor site among the three clinical matches and is selected by clinical/RNA annotation, without model advantage. Both AQR/PRPF8 events are displayed even though the two models fail AQR.

Existing count-head predictions give $\Delta_{v,t}^{\mathrm{ALT-REF}}=z_t(s_{v,\mathrm{ALT}},\Phi_{v,\mathrm{REF}})-z_t(s_{v,\mathrm{REF}},\Phi_{v,\mathrm{REF}})$ for Full; Sequence omits $\Phi$. The same reference conservation vector is used for both alleles. RNA inputs follow the source strand, whereas database matching uses genomic-plus alleles. All 13 source ASB calls favour REF, and the frozen source credits sum to seven for Full and ten for Sequence. No significance test, clinical AUROC, resampling or new inference was performed. Full/Sequence training histories remain incompletely matched. Figure~5e displays all four response arms as aligned horizontal bars, using $\Delta^{\mathrm{REF-ALT}}=-\Delta^{\mathrm{ALT-REF}}$ from the same frozen outputs. All four columns share the asinh axis (linear width 0.02, raw limits $[-0.12,0.88]$); no model or site has its own normalization. Filled positive bars agree with the source-preferred REF allele; open hatched negative bars favour ALT. Small group gaps separate the eight variants without aggregating any of the 13 event rows. Display IDs 1--8 follow the original frozen order; exact genomic/source identities are preserved in Figure~S32 and the accompanying CSV. Figure~S28g/h preserve the original asinh coordinates with linear width 0.02 and raw limits $[-1.2,1.2]$; ticks retain original log-count units. The flat- and permuted-$\Phi$ arms are saved responses at the same frozen Full weights; permuted responses average the three existing permutations and are not refitted models. The ELP1 zoom in Figure~S28g shows source RNA positions 42--58 (zero-based), with the focal donor +1 at position 50; its annotated exon--intron boundary is contextual annotation rather than a new splice assay. Figure~S28h shows source RNA positions 47--53 and each site's complete 101-nt reference-$\Phi$ vector on the common linear range $[-7,10]$. Source sequence and $\Phi$ identity across all events at each site were checked before drawing. One conservation row per variant spans its associated events; the 13 response rows retain every RBP/cell observation in frozen genomic order. Figure~S27b names every event.

\paragraph{Conservation and population context.} Center reference PhyloP agrees across every event at the same variant. Figure~S27c includes all 520 exact ClinVar variants by aggregate germline group, with medians and every unique-variant point; the three P/LP values do not support a general clinical-discrimination claim. The previous 241 allele-count-based 1000 Genomes frequencies remain unchanged in Data S25. Separately, current ClinVar \path{AF_TGP} is available for 428 of the 4,669 variants and is shown against reference PhyloP in Figure~S27d. The VCF does not specify that frequency's population version or denominator; neither is invented. An explicitly returned zero is retained, while an absent field remains missing. Among the eight displayed sites only ASL has a returned current TGP ALT frequency (0.00040). ALT frequency is not automatically minor-allele frequency, is not tumor prevalence and is not a binding or clinical gold label. Data S26 retains complete coverage/missingness, source-specific AF, current classifications, conditions and all joins.
